\documentclass{report}
\usepackage{settings}  % подключение настроек документа
\renewcommand{\labelenumi}{\arabic*.}

\begin{document}

%%%% Вставка необходимого титульника %%%%
%\include{TitlePages/diplomaTitle}  % Титульный лист дипломной работы 
\thispagestyle{titlepage}
%\newpage

    \begin{center}
    \linespread{1}
		Министерство науки и высшего образования  Российской Федерации\\
федеральное государственное бюджетное образовательное учреждение\\
высшего образования\\
«Иркутский государственный университет»\\
(ФГБОУ ВО «ИГУ»)\\
Факультет бизнес-коммуникаций и информатики 
    \end{center}
\vspace{1cm}    
\hspace{8cm} 
\begin{minipage}{0.5\textwidth}
\linespread{1}
%подписи для ВКР в курсовой можно удалить или заменить на настройку вертикального расстояния
  \begin{flushleft}
	\small{
    Кафедра естественнонаучных дисциплин \\%\vspace{-2cm}
Допускается к защите \\
и.о. зав. кафедрой, к.ф.-м.н., доцент \\
\rule{2,1cm}{0,1pt} А.Г. Балахчи\\		
            «\rule{0,9cm}{0,1pt}»\rule{2,7cm}{0,1pt} 202\,\rule{0,2cm}{0,1pt} г \\		}
		%\vspace{2еm}
		
		
		%\vspace{4.0em}
		\end{flushleft}
\end{minipage}


    \vspace{2cm} %вертикальное расстояние
    \begin{center}
   {\bf ВЫПУСКНАЯ КВАЛИФИКАЦИОННАЯ РАБОТА БАКАЛАВРА ИЛИ КУРСОВАЯ РАБОТА\\
    по направлению  \\
		09.03.03 Прикладная информатика\\
		%\vspace{1em}
		профиль\\
		«Разработка программного обеспечения»\\}
    \end{center}
		
		%\vspace{2.5em}%вертикальное расстояние
    \begin{center}
  		Название выпускной квалификационной работы или курсовой работы\\
    \end{center}
		
		\vspace{2.5em}
		
   \hspace{-5em} 


\noindent
\begin{minipage}[t]{0.5\textwidth}
  \begin{flushleft}
  \linespread{1}

		%\vspace{2em}
		
		
		%\vspace{4.0em}
		\end{flushleft}
\end{minipage}%\hspace{1em} 
\begin{minipage}[t]{0.5\textwidth}
  \begin{flushleft}
  \linespread{1}
	\small{
    Студент 3 курса
	очной формы обучения \\
  группа 14321\\
        \rule{1,8cm}{0,1pt} Воловиков Георгий Александрович\\
	\vspace{2em}
		
	Руководитель: уч. степень, уч. звание\\
            \rule{1,8cm}{0,1pt} Александр Викторович Кисилёв\\
\vspace{2em}
 Работа защищена:\\
        «\rule{0,9cm}{0,1pt}»\rule{2,7cm}{0,1pt} 202\,\rule{0,2cm}{0,1pt} г \\
с оценкой \rule{1,8cm}{0,1pt}\\
Протокол № \rule{0,9cm}{0,1pt}}
		%\vspace{2em}
		
		
		%\vspace{4.0em}
		\end{flushleft}
\end{minipage}

    \vfill
% Изменить год в надписи <<Иркутск 202*>> нужно в файле settings.sty на 221 строке
    
\newpage  % Титульный лист курсовой работы

\setcounter{page}{2} % начинаем нумерацию страниц
\tableofcontents  % это содержание, которое генерируется автоматически

\setcounter{chapter}{0} % установка счетчика глав
\setcounter{section}{0} % установка счетчика разделов
\setcounter{subsection}{0} % установка счетчика подразделов
\setcounter{equation}{0} % установка счетчика формул


\chapter*{ВВЕДЕНИЕ} 
\addcontentsline{toc}{chapter}{ВВЕДЕНИЕ} 

Нейросети сейчас повсюду --- их внедряют во все гаджеты, использует каждый человек, и они умеют решать огромное количество задач — от распознавания изображений до генерации текста и прогнозирования. При этом их работа, устройство и процесс создания для большинства людей до сих пор остаётся тайной. Это не случайно: разработкой нейросетей занимаются высококвалифицированные специалисты, которые проходят длительное обучение. Для создания собственной нейросети требуются продвинутые навыки программирования~(например, в Python или TensorFlow), математики~(линейная алгебра, статистика, оптимизация), а также теоретические знания о искусственном интеллекте. Эти требования создают высокий порог входа, который не позволяет широкой общественности понять и освоить процесс создания нейросетей. В частности, это закрывает такую возможность для следующих групп.

\begin{enumarabic}
\item Юных специалистов, которые хотят начать свой путь в Data Science. Они не могут полноценно изучать нейронные сети, поскольку нынешние инструменты требуют высоких навыков программирования, на освоение которых уходит несколько лет интенсивной практики.
\item Математиков, которые глубоко знают теоретическую основу нейросетей, но не изучают программирование напрямую и не владеют языками вроде Python.
\item Научных исследователей из нетехнических областей, которые хотят интегрировать нейросети в свои исследования, но не умеют программировать и сталкиваются с барьером реализации идей.
\end{enumarabic}

Отсутствие доступного программного обеспечения для разработки нейросетей без написания кода и необходимости глубоких знаний в математике ярко подтверждает {\bf актуальность} данной курсовой работы. Такие инструменты могли бы демократизировать доступ к технологиям ИИ, ускорить инновации и способствовать междисциплинарным исследованиям.

{\bf Объект исследования} --- процесс разработки нейронных сетей с использованием графических интерфейсов.

{\bf Предмет исследования} --- графические инструменты для разработки нейронных сетей.

{\bf Целью работы } является разработка MVP Desktop-приложения для графического создания нейронных сетей без необходимости в знаниях математики или программирования. Для достижения цели были поставлены следующие {\bf задачи}.

\begin{enumarabic}
\item Провести анализ существующих инструментов разработки нейронных сетей.
\item Разработать концепцию максимально гибкой и понятной графической разработки нейронных сетей.
\item Спроектировать архитектуру, обеспечивающую лучшую масштабируемость и производительность.
\item Разработать основу проекта — сам конструктор, механизм компиляции и графическую оболочку.
\item Апробировать продукт в реальных условиях.
\end{enumarabic}

{\bf Практическая значимость исследования} заключается в значительном снижении порога входа в разработку нейросетей и ускорении прототипирования проектов за счет реализации через графический интерфейс. Это также обеспечит наглядность всего процесса разработки, позволит популяризировать понимание работы нейросетей среди неспециалистов и облегчит обучение.

{\bf Теоретическая новизна исследования} заключается в том, что это первая платформа, которая позволяет максимально упростить процесс разработки и наглядно его продемонстрировать, при этом не ограничивая функциональность и позволяя проводить сложные эксперименты с продвинутыми архитектурами.

% обнуление счетчиков
\setcounter{section}{0} 
\setcounter{subsection}{0} 
\setcounter{equation}{0} 

\chapter{ТЕОРЕТИЧЕСКИЕ ОСНОВЫ И ОБЗОР СРЕДСТВ ВИЗУАЛЬНОГО ПРОЕКТИРОВАНИЯ НЕЙРОННЫХ СЕТЕЙ}

\section{Методология разработки нейросетей и требования к инструментарию}

Процесс разработки нейросети - это итеративный процесс постановки экспериментов и программирования их, множеством узкоспециализированных специалистов под каждый этап разработки. Разберем каждый этап процесса, представленного также на рисунке~\ref{fig:process}, подробнее.

\begin{enumarabic}
\item Сбор данных, осуществляется большими компаниями и требуют огромных мощностей и большого количества источников данных. Есть специальная профессия, которая настраивает работу сбора данных - Data Engineer. Считается, что в первую очередь от количества данных и зависит результирующее качество модели.
\item Обработка данных, играет не меньшую роль, чем сбор данных, потому что обучить на большом количестве данных очень дорого и часто можно получить тот же результат обучая на меньшем количестве данных более хорошего качества. В процессе сбора всегда возникают проблемы, шумят датчики, происходят непредвиденные события, солнечное затмения и т.д. Существует множество механик не только для очистки, но и для улучшения данных;
\item Разработка архитектуры - это итеративный процесс, который требует постоянных экспериментов и изменений. Обычно именно над этим занимаются разработчики большую часть времени и уделяют особое внимание, ведь именно тут больше всего требуются навыки программирования, математики и понимания нейронных сетей.
\item Обучение нейронной сети - самый долгий процесс. Хоть он и может показаться простым, но это вовсе не так. Стоит учитывать архитектуру данных, для определения наилучшего метода оптимизации. Рассчитывать лосс метрику исходят из архитектуры и задачи.
\item Валидация и тестирование - обучив модель, нужно ещё понять её качество, а для измерения качества используются разные метрики, выбор которых зависит от распределения данных, их объёма, задачи и многих других критериев~\cite{bib_goodfellow}.
\end{enumarabic}

\myfigure{0.9}{image1.png}{Процесс разработки нейросети}{process}

Несмотря на кажущуюся последовательность процесса разработки нейронной сети, на практике каждый его этап предъявляет собственные требования не только к знаниям разработчика, но и к используемым средствам проектирования. Именно поэтому из рассмотренного процесса разработки можно вывести требования к инструменту, который должен не просто предоставлять возможность создания модели, но и поддерживать разработчика на всех основных этапах работы с нейронной сетью.

\begin{enumarabic}
\item Простота в использовании. Разработка нейронных сетей уже сама по себе является сложной задачей, требующей понимания данных, архитектур моделей, обучения и оценки результатов. Поэтому инструмент разработки не должен дополнительно усложнять этот процесс за счёт громоздкого интерфейса, запутанной логики работы или избыточных действий со стороны пользователя.
\item Широкие функциональные возможности. Инструмент должен давать возможность не только собирать разные архитектуры нейросетей, но и решать другие задачи, которые возникают в процессе разработки. К ним относятся подготовка и обработка данных, настройка параметров модели, запуск обучения и другие вспомогательные действия. Чем больше этапов разработки поддерживает инструмент, тем удобнее использовать его в реальной работе.
\item Наглядность представления структуры нейронной сети. Архитектура модели является центральной частью, поэтому инструмент должен обеспечивать её удобное и понятное представление. Это упрощает понимание проектируемой системы, облегчает её анализ и делает процесс разработки более прозрачным. В нейронных сетях важно видеть не только то, какие слои используются, но и в каком порядке они стоят, как связаны между собой и какие параметры у них заданы. Даже если использовать одни и те же элементы, при разном их расположении можно получить совершенно разный результат.
\item Возможность настройки параметров модели. Разработка нейронной сети является итеративным процессом, в котором структура модели и параметры обучения многократно изменяются. Поэтому инструмент должен поддерживать удобное редактирование как архитектуры, так и основных параметров, влияющих на процесс обучения и итоговое качество модели.
\item Возможность интеграции в полноценные продукты. Результат разработки нейронной сети должен быть пригоден для дальнейшего использования вне среды проектирования. Это означает, что инструмент не должен быть полностью замкнутым на себе, а должен позволять переносить полученные решения в реальные программные системы, сервисы или исследовательские проекты.
\item Достаточная производительность при обучении моделей. Обучение нейронных сетей само по себе занимает много времени и требует больших вычислительных ресурсов. При этом каждая лишняя операция, замедление или неэффективное использование ресурсов увеличивает стоимость разработки. Поэтому инструмент должен обеспечивать достаточно высокую производительность, чтобы не замедлять процесс обучения и не делать работу с моделью дороже.
\end{enumarabic}

Подводя итог рассмотрению процесса разработки нейронных сетей, можно сделать вывод, что это многоэтапный, итеративный и ресурсоемкий цикл, каждый шаг которого традиционно требует высокой квалификации. Сформулированные выше требования показывают: чтобы сделать этот процесс доступным для начинающих специалистов и исследователей из нетехнических сфер, будущий инструмент должен не просто скрывать программный код, но и предоставлять гибкую, интуитивно понятную среду для управления архитектурой и параметрами модели.

Наиболее эффективным способом реализовать эти требования — в особенности обеспечить простоту использования и структурную наглядность "--- является применение парадигмы визуального программирования. В связи с этим, для выбора оптимального подхода к проектированию интерфейса и логики будущего конструктора, далее необходимо рассмотреть теорию графических языков программирования и проанализировать их применимость к задачам машинного обучения.


\section{Парадигма визуального программирования в контексте задач машинного обучения}

Визуальный язык программирования представляет собой способ создания программ, при котором логика системы задаётся не текстовым кодом, а с помощью визуальных элементов: блоков, узлов, пиктограмм, связей и диаграмм, размещаемых на рабочем поле по определённым правилам. В такой модели программа формируется через графическое конструирование и соединение компонентов, что делает структуру алгоритма более наглядной и облегчает восприятие логики~\cite{bib_vpl_theory}, особенно на этапах проектирования, обучения и прототипирования сложных систем.

Такие языки программирования обладают рядом преимуществ перед текстовыми. Представление логики в виде готовых компонентов, блоков и связей делает работу с программой более интуитивной, поскольку пользователю проще воспринимать графический список элементов и их назначение, чем текстовые описания функций и синтаксические конструкции. Кроме того, использование заранее определённых визуальных компонентов позволяет сократить объём ручного описания логики, благодаря чему разработка выполняется быстрее, а создаваемые программы нередко оказываются компактнее по структуре, чем на текстовых языках программирования.

Благодаря этим преимуществам графические языки программирования используются в тех сферах, где особенно важны наглядность и простота. В обучении они помогают формировать алгоритмическое мышление за благодаря наглядному представлению логики, для непрофильных специалистов служат средством работы с прикладными задачами без глубокого изучения кода и ускоряют прототипирования за счёт общего упрощения разработки.

\newpage При этом для понимания принципов работы графических языков программирования важно рассмотреть,из каких элементов они состоят.

\begin{enumarabic}
\item Синтаксис - это набор правил определяющий структуру. Он изучает, как слова объединяются в словосочетания и предложения, или как символы образуют правильные конструкции, обеспечивая корректное построение текста, речи или программы. В графических языка программирования он сохраняет свою роль и определяет как связываются графические единицы. Кроме того, часто это отображено графически: через цветы или формы соединений.
\item Семантика - это значение конструкций языка и описание того, что именно они делают при выполнении программы. Иначе говоря, если синтаксис определяет, как запись должна быть построена, то семантика определяет, какой смысл имеет эта запись и к какому результату приводит её использование. В графических языках программирования семантика может частично передаваться через визуальные средства представления, например через цвет, форму или расположение элементов. Благодаря этому пользователь может быстрее определить назначение блока и его роль в программе до детального изучения его содержимого.
\item Прагматика - это практическая сторона его использования, то есть то, насколько язык удобен, понятен и эффективен при решении реальных задач. Она связана с тем, насколько легко воспринимать программу, отслеживать её работу и получать результат в удобной для пользователя форме. В графических языках программирования прагматика часто проявляется через явность результата и непосредственную визуальную обратную связь: отображение выполнения работы программы через персонажей на сцене или визуализация методов и создания графиков.
\end{enumarabic}

Рассмотренные составляющие позволяют перейти от общего описания графических языков программирования к их классификации. Анализ видов графических языков программирования, на основе вышеперечисленных составляющих, позволит определить какой из них наиболее подходящий для нашей задачи.

\begin{enumarabic}
\item Блочные графические языки программирования — это языки, в которых программа составляется из готовых визуальных блоков, соединяемых между собой по правилам совместимости. В качестве примера можно привести Scratch, где синтаксис реализован через форму блоков и их сцепление по принципу пазла, семантика частично передаётся через цветовые категории блоков, а прагматика проявляется через сцену и спрайты, позволяющие сразу наблюдать результат выполнения программы. Для создания нейросетей данный вид подходит ограниченно, так как он удобен для обучения и описания простой логики, но хуже приспособлен к представлению сложных вычислительных графов.
\item Flowchart-подход является подвидом графовых языков программирования, в котором программа задаётся как блок-схема с последовательностью действий, ветвлений и циклов. Примером служит Flowgorithm, где синтаксис реализован через стандартные символы блок-схем и стрелки переходов, семантика определяется значением каждой фигуры, а прагматика состоит в пошаговой наглядности выполнения алгоритма. Такой подход полезен для изучения алгоритмов и построения последовательной логики, однако для создания нейросетей он малопригоден, поскольку ориентирован прежде всего на поток управления, а не на представление вычислительных узлов и передачи данных между ними.
\item DataFlow-подход также является подвидом графовых языков программирования, но в нём программа строится как схема передачи данных между узлами. Например Orange Data Mining, где синтаксис задаётся соединениями и виджетами, семантика определяется изображением на каждом узле, а прагматика проявляется в отображении результатов обработки в таблицах, графиках и других визуальных формах. Данный подход наиболее хорошо подходит для создания нейросетей, поскольку позволяет естественно представить слои, операции и поток данных внутри модели в виде направленного графа.
\item Языки форм или интерфейсно-ориентированные графические языки программирования — это языки, сочетающие визуальное программирование с проектированием интерфейса приложения. Примером является MIT App Inventor, где синтаксис задаётся через блоки и правила их соединения, семантика — через типы компонентов и событий, а прагматика усиливается тем, что пользователь одновременно видит форму приложения и его поведение. Этот подход в первую очередь ориентирован на создание приложений с интерфейсом, что не нужно для нейросетей.
\end{enumarabic}

\newpage Для более детального ознакомления с особенностями каждой категории, характерные примеры реализации различных видов визуальных языков программирования наглядно представлены ниже на рисунке~\ref{fig:vpl_types}.

\myfigure{0.9}{image2.png}{Виды графических языков программирования}{vpl_types}

Таким образом, проведенный анализ показывает, что визуальные языки программирования являются эффективным средством снижения порога входа в сложную предметную область. Для задачи конструирования нейронных сетей наиболее перспективным является подход DataFlow, поскольку он максимально естественно отражает архитектуру искусственных нейронных сетей - в виде направленного графа узлов и потоков данных.

Успешная реализация такого инструмента требует тщательного проектирования его структурных компонентов: синтаксис должен предотвращать создание некорректных архитектур, семантика - обеспечивать прозрачность работы каждого слоя, а прагматика - давать мгновенную обратную связь в процессе обучения. Именно учет этих составляющих позволит создать не просто визуальную оболочку, а полноценную образовательную и исследовательскую среду. Для того чтобы определить, какие архитектурные и интерфейсные решения наиболее эффективны на практике, а каких ошибок следует избегать, далее необходимо провести сравнительный анализ существующих визуальных инструментов разработки ML-моделей.


\section{Сравнительный анализ существующих платформ для разработки исскуственного интеллекта}

Для успешного проектирования собственного визуального конструктора нейронных сетей недостаточно опираться только на теоретическую базу, необходимо также изучить практический опыт реализации подобных систем. Это позволит, с одной стороны, выявить наиболее удачные интерфейсные и функциональные решения, чтобы взять их за основу и сформировать примерный образ будущего продукта, а с другой, обнаружить критические недостатки и ограничения аналогов, чтобы избежать их и предложить более совершенный подход в разрабатываемом приложении.

Оценка существующих платформ будет на основе требований, выведенных в предыдущих главах: простота, функциональные возможности, наглядность, тонкая настройка, возможность интеграции, производительность и вид инструмента.

Существующие аналоги можно предварительно разделить на 4 крупные группы.

\begin{enummarker}
\item Фреймворки. Это полноценные библиотеки, дающие полные возможности к реализации своих нейросетей. Они позволяют на уровне отдельных слоёв разработать модель, что обеспечивает высокую гибкость и производительность. При этом, они сложны в освоении, потому что нужно изучить не только их, но и предварительно язык программирования. Не смотря на то, что многие фреймворки предлагают возможность визуализации архитектуры, наглядность всего решения в целом, всё равно отсутствует. При этом для полноценных решений в больших компаниях, именно эти варианты подходят лучше всего.
\item Сервисы работающие по принципу чёрной коробки. Принцип чёрной коробки заключается в том, что нам не нужно знать, как работает чёрная коробка и что находиться внутри её, важен лишь результат и как мы можем на него воздействовать. Исходя из этого принципа данные сервисы слишком упрощают разработку нейронных сетей, не позволяя взглянуть внутрь. Такой принцип также не подразумевает под собой большое количество возможностей. Чаще всего они работают по методу загрузить данные и получить модель. Хотя такой метод и очень быстрый для разработки, а также удобный для интеграции в другие системы, но совершенно не подходит для объяснений внутренней составляющей нейронных сетей, из-за полного отсутствия наглядности.
\item Сервисы визуализации работы нейронных сетей. Есть область исследований, которая нацелена на интерпретацию работы нейронных сетей. Эти инструменты из этой области - их главная задача, не позволить создать нейронные сети и не объяснить, как они работают, а показать, как они думают. Это очень полезно, для глубокого понимания работы моделей и их прогнозирования их дальнейшего поведения. Их предназначение в первую очередь в интерпретации нейронных сетей, а не в их создании, из-за которых такие инструменты имеют малый функционал.
\item DataFlow-платформы для ML. Отдельной группой я решил вынести DataFlow инструменты, которые реализуют разработку моделей машинного обучения. Суть их в том, что каждый узел реализует какую-то свою логику - загрузка данных, создание модели, расчет метрик или отрисовка графиков. Такой подход возможен, благодаря простой работе с моделями машинного обучения, вероятнее всего, именно из-за этого они не рассматривают реализации нейронных сетей тем же методом, что является их главным недостатком.
\end{enummarker}

Опираясь на выявленные преимущества и ограничения рассмотренных групп, можно сделать вывод, что первые три категории решений не в полной мере удовлетворяют заявленным требованиям к будущему продукту. Кодовые фреймворки сохраняют слишком высокий порог входа, no-code сервисы работают по принципу «чёрного ящика» и лишены образовательной наглядности, а демонстрационные утилиты обладают критически ограниченным функционалом.

В то же время, как было обосновано в предыдущем разделе, именно парадигма DataFlow наиболее естественно отражает архитектуру искусственных нейронных сетей. Следовательно, для формирования детального образа разрабатываемого приложения и поиска оптимальных интерфейсных решений наибольший научно-практический интерес представляют платформы из четвертой группы — DataFlow-инструменты для машинного обучения. Тщательное изучение их функциональных особенностей позволит перенять лучшие практики и найти способы решения их текущих недостатков.

Для этого далее рассмотрим DataFlow платформы более подробно.

Orange Data Mining — это мощная открытая~(open-source) платформа для визуального программирования, анализа данных и машинного обучения. Инструмент представляет собой классический пример реализации парадигмы DataFlow, где весь процесс разработки выстраивается на интерактивном холсте, что изображено на рисунке~\ref{fig:orange_mining}~\cite{bib_orange}.

\myfigure{0.9}{image3.png}{Интерфейс Orange Data Mining}{orange_mining}

Принцип работы Orange Data Mining, такой же, как и у всех DataFlow инструментов. Мы размещаем на рабочем пространстве узлы, каждый из которых отвечает за своё действие, а потом соединяем их, тем самым создавая программу. Семантика тут очень хорошо показана, тем что, каждый узел это своя картинка, по которой интуитивно понятно, что он выполняет. Прагматика делается через графики и визуализации методов. Но синтаксис реализован плохо, все связи визуального одного вида и цветов, хотя на самом деле это не так и платформа ограничивает пользователя в некорректных связях.  

Благодаря DataFlow структуре мы видим весь процесс разработки в виде единой структуры - графа, этим обеспечивается полная наглядность и простота в использовании. Но инструмент не реализует более сложные варианты разработки. У него нет узлов для создания полноценных, сложных нейронных сетей, только простейшие перцептроны, а также он не предоставляет возможности для тонкой настройки каждого из узлов и не даёт возможности для выгрузки итоговой модели, для дальнейшего использования в других проектах, а производительность оставляет желать лучшего.

Не смотря на это, изначально платформа создавалась в академической среде, поэтому она имеет сильный уклон в сторону визуализации данных и статистики. Этот инструмент лучше всего подойдет студентам, преподавателям, а также аналитикам и исследователям из нетехнических сфер, где не нужно сильное углубление или варианты сложной разработки.

SMILE~(Smart Machine Learning Environment) — это учебно-исследовательская платформа визуального проектирования моделей машинного обучения, разрабатываемая в университете ИТМО. Инструмент создан специально для образовательных нужд, чтобы продемонстрировать полный цикл создания ML-моделей в виде наглядного конвейера. Основная целевая аудитория платформы - школьники старших классов, студенты младших курсов и участники профильных олимпиад, которым нужно освоить концепции машинного обучения без погружения в сложный синтаксис языков программирования.

\myfigure{0.9}{image4.png}{Интерфейс платформы SMILE}{smile_platform}

Система работает на основе блочного конструктора в браузере, где логика обработки данных и построения модели собирается в виде графа. Главная техническая особенность SMILE заключается в том, что после сборки визуального графа система автоматически транслирует эти блоки в исполняемый Python-код. Затем этот код отправляется на сервер, где происходит компиляция, обучение модели на вычислительных мощностях университета и возврат метрик обратно в графический интерфейс пользователя. Эта особенность является одновременно главным преимуществом и недостатком платформы. С одной стороны, мы можем получить сразу Python код, а дальше работать уже с ним. С другой стороны, реализация через веб, накладывает большие вычислительные ограничения, из-за чего сделать такую платформу общедоступной для всех не представляется возможным.

По критериям простоты и наглядности SMILE справляется со своей образовательной задачей. При этом, там опять визуально отсутствует синтаксис и даже больше, нет никаких ограничений по связям узлов, что позволяет пользователю создавать ошибки. Семантика хоть и задаётся через различные цвета блоков, но на самом деле цветов очень мало, из-за чего визуально сложно определить за что они отвечают. В качестве прагматики предлагается результирующий код программы.. Разработкой занимается небольшая команда университета, для которой широкий функционал не приоритетен. При этом конструктор очень производителен, потому что не несёт за собой никаких вычислительных расходов, а просто создаёт Python код, по этой же причине, его модели очень просто интегрировать в уже существующие проекты.

Deep Learning Studio — специализированная визуальная платформа для разработки нейросетей. В отличие от аналогов, она полностью сфокусирована на глубоком обучении, позволяя быстро прототипировать архитектуры. Система доступна в облачном и Desktop-форматах.

\myfigure{0.9}{image5.png}{Интерфейс Deep Learning Studio}{dls_platform}

Интерфейс системы реализован на 2-ух уровнях. Тут есть несколько вкладок, каждая из которых отвечает за свой этап разработки: подготовка данных, разработка архитектуры, подбор гипер параметров, обучение и визуализация результатов. Нас в первую очередь интересует вкладка для разработки архитектуры. Тут представлен такой же DataFlow интерфейс, как в предыдущих работах. Не смотря на то, что решение разложить процесс разработки по вкладкам кажется логичным, на самом деле это очень сильно усложняет интерфейс. Пользователю нужно уметь работать не только с DataFlow интерфейсом разработки самой архитектуры, но и уметь работать с табличным интерфейсом подготовки данных, интерфейсом подобра гипер параметров и т.д.

При этом интерфейс всё равно остаётся гораздо проще профессиональных фреймворков. Благодаря DataFlow интерфейсу, архитектура нейросети тоже остаётся очень наглядной, но остальной процесс разработки - нет. Из-за того, что мы собираем тут только архитектуру нейросети, нам не нужен синтаксис, потому что все блоки могут быть соединены в произвольном порядке. Семантика при этом тут остаётся очень наглядной, благодаря изображениям на каждом узле. 

Главной технической особенностью платформы, является то, что она базируется на TensorFlow. Это даёт ей широкие функциональные возможности~\cite{bib_tensorflow}, быструю производительность, а также возможность интегрировать готовые модели в любые проекты с этим фреймворком. Разработчики также подумали и о тонкой настройке параметров каждого узла, что происходит в отдельном окне.

MATLAB Deep Learning Toolbox - это профессиональная проприетарная среда для проектирования, обучения и анализа нейронных сетей от компании MathWorks. В отличие от легковесных образовательных платформ, это часть огромной экосистемы MATLAB, исторически ориентированной на сложные инженерные и математические расчеты. Основная целевая аудитория данного инструмента - инженеры, научные сотрудники, специалисты по обработке сигналов и компьютерному зрению, а также студенты профильных технических вузов, которым необходимо интегрировать алгоритмы глубокого обучения в более сложные системы моделирования внутри MATLAB.

Процесс разработки очень похож на Deep Learning Studio. Тут также нет одного интерфейса, который бы реализовывал сразу весь процесс разработки, предлагается только разработка архитектуры нейросети через DataFlow интерфейс. Уникальной особенностью платформы является мощный встроенный инструмент статического анализа сети, который до нажатия кнопки «Обучить» проверяет архитектуру на наличие синтаксических ошибок, пропущенных связей и строит детальную таблицу несоответствия размерностей тензоров между слоями.

\myfigure{0.9}{image6.png}{Интерфейс MATLAB Deep Learning Toolbox}{matlab_toolbox}

Наглядность получается двоякой, с одной стороны мы отлично видим архитектуру нейросети, но остальные этапы разработки нет. При этом порог входа остается сравнительно высоким, потому что для обработки данных и результатов требуется знание сложной системы MATLAB. Функциональные возможности сравнительно высокие, уже сейчас предлагается список из достаточно большого количества узлов, а дополняя это возможностями самого MATLAB получаются совсем неограниченные возможности Разработчики также предусмотрели настройку каждого узла по отдельности. Производительность невероятная высокая, не смотря на то, что проект не основан на каком-то уже готовом фреймворке. Также есть возможность экспорта в ONNX формат. Это идеальный вариант для интеграции модели в производственные системы.

Проведенный сравнительный анализ существующих платформ позволил выделить их ключевые достоинства, которые целесообразно интегрировать в единую систему, а также выявить критические недостатки, требующие устранения. На основе этого синтеза формируется концептуальный облик разрабатываемого конструктора, который должен стать сбалансированным инструментом, сочетающим интуитивность образовательных сервисов и вычислительную мощность профессиональных сред.

Простота и наглядность проектируемого приложения будут опираться на успешный опыт Orange Data Mining. Главная идея заключается в реализации всего процесса разработки нейронной сети, от загрузки и обработки данных, до создания архитектуры и её валидации, в рамках единого DataFlow-интерфейса на одном рабочем холсте. Это обеспечит полную наглядность, позволяя пользователю видеть всю логику проекта как целостный связный граф вычислений без необходимости переключаться между разными окнами.

Функциональные возможности и высокая производительность системы будут реализованы по аналогии с Deep Learning Studio. Для этого проект будет базироваться на современной и высокопроизводительной библиотеке глубокого обучения. Такой подход позволит реализовывать сложные решения, которые будут быстро и надёжно работать, а также быстро масштабировать количество узлов.

В вопросах возможности интеграции за эталон принимается подход MATLAB Deep Learning Toolbox. Разрабатываемый конструктор должен поддерживать экспорт готовой и обученной модели в общий универсальный формат, что обеспечит возможность её дальнейшего свободного использования в реальных программных продуктах и исследовательских проектах вне среды конструктора.

Тонкую настройку никто не реализовал в идеальном виде. Не смотря на то, что она есть у нескольких платформ, но все они создают её в отдельном окне, что усложняет интерфейс. Вместо этого лучше будет реализовать её прямо внутри узлов, но убирая лишние параметры, для того, чтоб не перегружать интерфейс.

Особое внимание при проектировании будет уделено решению проблем, выявленных у конкурентов, через глубокую проработку структурных компонентов визуального языка.

\begin{enummarker}
\item Синтаксис. В отличии от остальных аналогов, в новом инструменте будет цветовая индикация коннекторов, для интуитивно понятного, отображения возможных связей.
\item Семантика. Не смотря на то, что реализация семантики через изображение, как у Deep Learning Studio и других, будет более интерпретируемой, я всё же решил сделать как у SMILES. Создание своих узнаваемых изображений на каждый логическо разных вид узлов будет накладывать большие затраты на их разработку. В свою очередь, простое цветовое разделение по группам даст возможность быстрее внедрять новые функциональные узлы.
\item Прагматика. Обратная связь и практическая наглядность будут реализованы в стиле Orange Data Mining, где результаты работы отражаются в графиках.
\end{enummarker}

Таким образом, разрабатываемый MVP конструктора призван занять пустующую нишу инструментов, которые предлагают профессиональную глубину и производительность ML-разработки при сохранении максимальной прозрачности визуального языка и низкого порога входа.


\section{Обоснование выбора технологического стека}

После формирования концептуального облика разрабатываемого конструктора и определения функциональных требований к нему, необходимо перейти к выбору технологического стека, который позволит реализовать MVP. Выбор правильных инструментов критически важен, так как от них зависит производительность, масштабируемость и сложность разработки платформы.

Базовым языком программирования для реализации всего проекта был выбран Python. Данный выбор обусловлен несколькими факторами. Во-первых, Python является абсолютным мировым стандартом в сфере Data Science~\cite{bib_python_docs} и машинного обучения, обладая самой обширной экосистемой библиотек. Во-вторых, он отличается лаконичным синтаксисом и простотой разработки, что позволяет значительно ускорить процесс создания прототипов и реализации сложной архитектуры программного обеспечения.

Ранее мы уже обозначали, что для реализации продукта с максимальной масштабируемостью, надёжностью и производительностью, стоит его реализовывать на основе уже готового фреймворка для искусственного интеллекта. На сегодняшний день, наиболее распространены и совершены следующие 3 фреймворка.

\begin{enummarker}
\item PyTorch - самый популярный фреймворк на сегодняшний день. Главной его особенностью является мощный динамический граф вычислений, а также самые широкие функциональные возможности без пере усложнений. При этом он остаётся всё ещё слишком сложным для создания проекта на его основе. Процесс обучения для каждой нейросети в нём реализуется по разному, а сама нейросеть собирается только в виде класса.
\item JAX - высокопроизводительная библиотека от Google, предназначенная для сложных математических вычислений с использованием автоматического дифференцирования и XLA-компиляции. JAX отлично подходит для передовых научных исследований, но является слишком низкоуровневым и сложным инструментом. Использование JAX для базового MVP конструктора нейросетей является избыточным.
\item TensorFlow~(с API Keras) — зрелая и мощная платформа от Google, ориентированная как на исследования, так и на внедрение в готовые продукты. Ключевым преимуществом TensorFlow для разрабатываемого проекта является наличие Functional API~\cite{bib_keras_api}. Данный подход позволяет представлять слои нейронной сети как узлы направленного графа, что идеально ложится на выбранную нами концепцию DataFlow. Кроме того, TensorFlow имеет встроенные высокоуровневые методы обучения и валидации, которые легко инкапсулировать внутрь визуальных блоков, избавив пользователя от создания цикла обучения вручную.
\end{enummarker}

Определившись с ядром проекта - фреймворком для нейросетей, можно приступить и к выбору графического фреймворка. Поскольку разрабатываемый продукт является Desktop-приложением, графическая библиотека должна удовлетворять главному требованию — быть кроссплатформенной, чтобы конструктор мог одинаково корректно работать на операционных системах Windows, macOS и Linux. Рассмотрим основные графические кроссплатформенные фреймворки на Python.

\begin{enummarker}
\item PyQt~(и его аналог PySide) - индустриальный стандарт создания GUI на Python. Это невероятно мощный инструмент, для которого существуют сторонние библиотеки графовых узлов. Тем не менее, PyQt является очень «тяжелым» фреймворком. Он несет за собой массивный объем зависимостей, требует сложной настройки окружения и потребляет значительное количество оперативной памяти, что может замедлить работу приложения, которое и так будет нагружено вычислениями нейронных сетей.
\item Flet и Kivy - современные фреймворки, позволяющие создавать красивые и кросс-платформенные интерфейсы. Однако их существенным недостатком для данного проекта является отсутствие встроенных механизмов для создания графовых редакторов~(Node Editor). Реализация холста с узлами, портами и изогнутыми линиями связей на базе этих библиотек потребовала бы написания огромного количества низкоуровневого кода для отрисовки графики, что нецелесообразно на этапе создания MVP.
\item DearPyGUI - это современный, минималистичный графический фреймворк, использующий рендеринг на GPU. Он отличается высокой скоростью работы, малым весом и полным отсутствием зависимостей. Решающим преимуществом DPG стало наличие встроенной подсистемы для создания редакторов графа~\cite{bib_dearpygui}.
\end{enummarker}

Подводя итог выбору инструментария, можно утверждать, что связка из языка Python, фреймворка машинного обучения TensorFlow и графической библиотеки DearPyGUI образует оптимальный технологический стек для разрабатываемого MVP. TensorFlow обеспечивает надежное и функциональное вычислительное ядро, идеально совместимое с графовыми архитектурами, а DearPyGUI предоставляет готовый, легковесный, аппаратно-ускоренный интерфейс для их визуального конструирования. Выбор данных инструментов полностью отвечает сформированным ранее требованиям к продукту и закладывает прочный фундамент для перехода к практической части работы "--- проектированию программной архитектуры конструктора.


\section*{Выводы по главе}
\addcontentsline{toc}{section}{Выводы по главе}
Был проведен подробный анализ процесса разработки нейронных сетей и сформулированы требования к современному инструменту проектирования моделей машинного обучения. На основе изучения теории графических языков программирования был обоснован выбор парадигмы DataFlow как наиболее естественного способа представления архитектур нейронных сетей в виде направленного ациклического графа.

Сравнительный анализ существующих аналогов, таких как Orange Data Mining, SMILE, Deep Learning Studio и MATLAB, выявил отсутствие на рынке инструментов, которые бы сочетали в себе профессиональную глубину разработки с высокой наглядностью и низким порогом входа. По итогам обзора инструментария для реализации MVP-проекта был выбран технологический стек, включающий язык Python, вычислительное ядро TensorFlow и графический фреймворк DearPyGUI. Данные решения позволяют создать производительную кроссплатформенную среду, способную стать базой для эффективного обучения и прототипирования нейросетевых моделей.


\chapter{ПРОЕКТИРОВАНИЕ, ПРОГРАММНАЯ РЕАЛИЗАЦИЯ И ОЦЕНКА ЭФФЕКТИВНОСТИ КОНСТРУКТОРА}
\setcounter{section}{0}
\setcounter{subsection}{0}
\setcounter{equation}{0}

\section{Разработка программной архитектуры}

Основываясь на принципах, сформулированных в теоретической части работы, и концепции атомарности операций, была разработана архитектурная модель программного комплекса. В её основу положен тезис о том, что каждый узел графа является инкапсулированной оболочкой для конкретной Python-функции. В рамках данной модели связи между узлами определяют не только визуальную последовательность, но и поток данных~(DataFlow), где результат выполнения функции одного узла становится входным аргументом для последующего. Подобный подход обеспечивает максимальную гибкость системы, позволяя реализовывать сложные алгоритмические структуры, и гарантирует высокую масштабируемость за счет построения узлов на базе готовых библиотечных функций.

Программную архитектуру можно разделить на два ключевых функциональных модуля: механизм генерации узлов и систему управления динамическим вычислительным графом.

Рассмотрим подробнее реализацию первого модуля. Интерпретируя узел как совокупность аргументов, передаваемых в функцию, становится возможным автоматизировать формирование интерфейса на основе сигнатуры этой функции. В языке Python состав и типы аргументов определяются через механизм аннотаций. Несмотря на то, что использование аннотаций является опциональным, их применение считается стандартом при разработке качественного ПО. Интроспекция существующих аннотаций в функциях TensorFlow могла бы значительно ускорить процесс разработки.

Однако в ходе реализации было выявлено отсутствие необходимых аннотаций в ряде методов TensorFlow. Для решения данной проблемы был разработан специализированный конфигурационный файл, в котором фиксируются функции, их параметры, аннотации, идентификаторы узлов и их типы. Создание подобного конфигурационного слоя обеспечило не только масштабируемость, но и гибкость настройки: стала возможной реализация интуитивно понятных наименований узлов и интеграция документации.

Процесс создания каждого узла делегирован специализированному менеджеру — node\_builder. Данный модуль извлекает конфигурацию узла из файла и инициализирует объект, формируя на основе аннотаций соответствующие поля ввода. Также на данном этапе интегрируется справочная информация и определяются входные/выходные порты для установления связей в графе.

\myfigure{0.9}{image7.png}{Создание узла}{arch1_img}

Предложенная архитектура обеспечила масштабируемость по количеству доступных узлов, однако в ходе развития проекта возникла проблема масштабируемости самих типов аннотаций. Рост числа стандартных типов и усложнение их логики привели к необходимости создания комплексных полей ввода, состоящих из нескольких типов данных~(например, для параметров сверточных слоев). Реализация таких элементов требовала проверки не только типов, но и сложных объектов. Кроме того, возникла потребность в специфическом типе ввода, требующем доступа к объектам других узлов. Сложность заключалась в том, что информация о связях в библиотеке DearPyGUI~(DPG) хранится в атрибутах — специальных структурах узла, обеспечивающих визуальное и логическое соединение, — что требовало выхода на более высокий уровень абстракции.

\myfigure{0.4}{image8.png}{Проблема типов аннотаций}{arch2_img}

Для минимизации сложности кода на данном этапе был применен шаблон проектирования «Фабрика»~\cite{bib_gamma_patterns}. Система была организована как совокупность нескольких фабрик, интегрированных в единый интерфейс. Предполагалось, что каждая фабрика будет отвечать за создание полей ввода на основе определенных принципов: либо на базе типов данных, либо на основе конкретных объектов.

\myfigure{0.7}{image9.png}{Использование фабрик}{arch3_img}

Несмотря на сокращение дублирования кода, временные затраты на поддержку системы возросли из-за избыточной сложности взаимодействия между фабриками. Процесс добавления нового функционала требовал от разработчика детального анализа всей цепочки наследования и взаимодействия пяти различных модулей.

В целях упрощения системы было принято решение отказаться от использования фабрик в пользу концепции диспетчеризации. Диспетчеризация предполагает вызов различных реализаций функции в зависимости от типов и значений передаваемых аргументов. Данный подход расширяет принципы полиморфизма и перегрузки. В языке Python для этих целей использовался стандартный инструмент singledispatch, дополненный авторской реализацией декоратора factorymethod. Это позволило выстроить прозрачную логику ветвления, представленную на рисунке~\ref{fig:arch4_img}.

\myfigure{0.7}{image10.png}{Логика ветвления~(система диспетчеризации)}{arch4_img}

Данная реализация упростила понимание кодовой базы, сохранив гибкость, присущую фабрикам. В частности, это позволило внедрить динамические выпадающие списки на основе перечислений. Однако возникла необходимость расширения функционала типов: требовалось не только считывать данные из полей ввода, но и программно изменять их значения для повышения прагматики системы~(например, для отображения результатов расчетов и графиков непосредственно в узлах при компиляции).

Окончательным архитектурным решением стало проектирование собственной иерархии типов аннотаций. Каждый тип инкапсулирует логику создания визуального компонента, извлечения данных и их записи. В новой структуре параметры узла описываются через объект Parameter, который инициирует создание атрибута, а его поведение определяется классом Annotation~(например, AInteger, AFloat, ANode, AFile, AEnum и др.). Каждый класс реализует три обязательных метода: отрисовка интерфейса, получение данных и вставка данных. Таким образом был достигнут полный контроль над типизацией аргументов, упрощена кодовая база и обеспечена высокая расширяемость.  

\myfigure{0.9}{image11.png}{Иерархия аннотаций}{arch5_img}

Реализованная архитектура узлов решает задачу масштабируемости интерфейса, в то время как механизм компиляции определяет логику взаимодействия компонентов. Алгоритм компиляции графа включает четыре аспекта: модель хранения графа, логику компиляции отдельного узла, определение точек входа и порядок обхода.

Для хранения структуры графа была выбрана объектно-ориентированная модель, где граф представлен совокупностью связанных классов. Был разработан базовый класс AbstractNode, инкапсулирующий методы управления~(установление/удаление связей, инициализация) и внутренний алгоритм компиляции. Метод компиляции извлекает значения через механизмы аннотаций и передает их во внутреннюю функцию, сохраняя результат для последующей передачи по графу. Класс AbstractNode имеет специализированных наследников, что обеспечивает строгую типизацию соединений и предотвращает построение некорректных архитектур.

Поскольку концепция DataFlow-платформы базируется на направленном ациклическом графе~(DAG), для корректной работы необходимо точно определять начальные узлы — те, которые не имеют входящих связей и могут быть вычислены независимо. Первоначальное решение, основанное на алгоритме поиска в ширину~(BFS) от произвольного узла, обладало высокой асимптотической сложностью, требуя полного обхода графа.

Для оптимизации был разработан механизм кэширования начальных узлов в режиме реального времени. Данное решение опирается на событийную модель редактирования графа.

\begin{enumarabic}
\item Все вновь создаваемые узлы инициализируются как начальные.
\item При установке связи, если у узла появляется родительский элемент, он исключается из списка начальных.
\item При разрыве связи, если узел теряет последнего родителя, он вновь регистрируется как начальный.
\item Перед удалением узла происходит предварительная деактивация всех его связей.
\end{enumarabic}

Подобный событийный подход позволил полностью исключить фазу предварительного поиска точек входа, что значительно ускорило процесс компиляции и сделало систему более отзывчивой.

Таким образом, спроектированная программная архитектура, основанная на принципах атомарности операций и парадигме DataFlow, обеспечивает необходимую гибкость и масштабируемость системы. Переход от простейшей реализации на условиях и избыточных шаблонов проектирования к иерархической структуре типов аннотаций позволил не только упростить кодовую базу, но и достичь полного контроля над процессом генерации интерфейса узлов. Реализованная объектно-ориентированная модель графа с базовым классом AbstractNode гарантирует корректность построения архитектур нейронных сетей, а внедрение механизма событийного кэширования начальных узлов оптимизировало процесс компиляции направленного ациклического графа. В совокупности данные решения формируют надежный фундамент для дальнейшей реализации пользовательского интерфейса и функциональных возможностей конструктора.


\section{Проектирование графического интерфейса}

На основании теоретического анализа, проведенного в первой главе, в качестве концепции интерфейса была выбрана парадигма DataFlow. Данная концепция обычно реализуется посредством механизма Drag-and-Drop между двумя окнами: в левом окне представлен перечень доступных блоков, а в правом расположена рабочая область, в которой осуществляется конструирование вычислительного графа путем перетаскивания и связывания узлов. При проектировании было критически важно реализовать три ключевых компонента визуального языка программирования: синтаксис, семантику и прагматику. Соблюдение данных принципов позволяет создать интуитивно понятный, гибкий и эффективный в освоении и использовании конструктор.

Вследствие технических ограничений библиотеки DearPyGUI, исключающей инициализацию элементов типа node вне контейнера node\_editor, на текущем этапе узлы в левой панели отображаются в виде кнопок. Все доступные узлы распределены по функциональным группам внутри выпадающих списков. В перспективе планируется реализация механизма имитации узлов для оптимизации пользовательского опыта.

Рабочая область реализована с использованием встроенных инструментов DearPyGUI, однако следует выделить несколько специфических особенностей ее функционирования. Во-первых, в рабочей области по умолчанию инициализируется узел входных данных нейросети. Данный элемент является точкой входа, обеспечивает структурную целостность графа и служит базой для определения координат курсора в рабочем пространстве. Во-вторых, графический фреймворк предоставляет исключительно визуальные компоненты, что обусловило необходимость разработки логики управления графом с нуля: алгоритмов создания, связывания, развязывания, размещения и удаления узлов.

\myfigure{0.9}{image12.png}{Рабочая область интерфейса}{workspace_img}

Этого функционала достаточно для основной части графического интерфейса, однако для создания интуитивно понятной среды, пригодной для образовательной деятельности, была проведена детальная проработка трех компонентов визуального языка программирования.

\begin{enumarabic}
\item Синтаксис. В рамках конструктора синтаксис определяет правила взаимодействия блоков. Реализована цветовая индикация коннекторов: каждому типу данных соответствует определенный цвет. Программная логика блокирует установку связей между портами разного цвета, что визуально информирует пользователя о совместимости узлов.
\item Семантика. Отвечает за функциональное содержание узлов. Реализована цветовая дифференциация групп: слои нейронных сетей маркируются красным цветом, этапы обучения — фиолетовым, модули обработки данных — бирюзовым. Наличие параметров для настройки непосредственно внутри узлов позволяет пользователю однозначно интерпретировать назначение конкретного элемента.
\item Прагматика. Реализуется через инструменты оперативной визуализации данных. Был разработан специализированный тип аннотаций — AFigure, обеспечивающий интеграцию графиков библиотеки Matplotlib непосредственно в узлы конструктора. Это позволяет пользователю анализировать результаты вычислений и метрики обучения в режиме реального времени прямо на рабочем поле.
\end{enumarabic}

\myfigure{0.9}{image13.png}{Реализация прагматики и графиков}{pragmatics_img}

В процессе реализации интерфейса также был расширен функционал DearPyGUI для соответствия архитектурным требованиям. Был создан менеджер тем. Изначально DearPyGUI предполагает описание стилей непосредственно в коде с их последующим наложением. Подобный подход усложнил бы создание запланированного дизайна и ограничил бы масштабируемость системы. Менеджер тем оперирует конфигурационными файлами формата JSON, в которых прописаны все графические параметры. Ключевой особенностью менеджера является поддержка наследования тем, реализованная на основе логического объединения словарей Python. Таким образом, можно накладывать новые стили без сброса предыдущих настроек. Благодаря такому подходу возможна оперативная разработка новых визуальных тем без изменения исходного кода проекта.

Также была расширена механика функций обратного вызова~(callbacks). DearPyGUI содержит множество обработчиков событий~(нажатие кнопок, движение мыши и т. д.), однако изначально фреймворк поддерживает назначение лишь одной функции на каждое действие. Это создает трудности при необходимости использования разных обработчиков для различных графических элементов. В связи с этим был реализован менеджер обратного вызова, который аккумулирует все функции на каждое действие и определяет принадлежность события к конкретному элементу. При активации события менеджер производит последовательный вызов всех привязанных к нему функций.

Разработанная оболочка отвечает всем требованиям: синтаксическая корректность обеспечена типизированными коннекторами, семантическая — цветовой дифференциацией групп, а прагматика — инкапсуляцией настроек и выводом вычислений внутри узлов. Мини-карта, иерархическая панель и система менеджеров делают конструктор завершенным масштабируемым инструментом для эффективного проектирования сложных нейросетей в едином пространстве.

\myfigure{0.9}{image14.png}{Полный графический интерфейс}{interface_full_img}

\section{Верификация системы и автоматизированное тестирование}

Для обеспечения надежности разработанного MVP-конструктора и проверки корректности работы механизма компиляции графа был проведен комплекс мероприятий по тестированию. Процесс верификации включал в себя несколько уровней: ручное функциональное тестирование, модульное тестирование бизнес-логики, специализированное UI-тестирование и интеграционные тесты.

    В первую очередь проведено ручное функциональное тестирование продукта. На этом этапе оценивалась корректность работы алгоритмов и эргономичность интерфейса. К апробации привлекли профильных специалистов и целевую аудиторию — школьников. Благодаря их обратной связи был выявлен и устранен ряд визуальных и логических дефектов. Примеры исправленных в ходе тестирования багов представлены на рисунке~\ref{fig:bugs_img}.

\myfigure{0.9}{image15.png}{Устраненные баги}{bugs_img}

Параллельно с разработкой функционал проекта покрывался автоматизированными тестами. Первостепенное внимание уделялось логической части приложения, изолированной от графического интерфейса. Проверялись функции внутри вычислительных узлов, алгоритмы связывания элементов внутреннего графа, а также работа системных менеджеров. База автотестов регулярно дополнялась регрессионными сценариями на основе багов, выявленных при ручном использовании системы.

Особой задачей стало специализированное UI-тестирование компонентов, тесно связанных с графическим интерфейсом. Для них был разработан класс тестового окружения, который инициализировал графический контекст перед запуском тестов для корректной работы фреймворка DearPyGUI. Важной особенностью данного подхода стало то, что графические элементы создавались и функционировали корректно, но физически не отрисовывались на экране. Это позволило значительно ускорить процесс выполнения тестов при полном сохранении их достоверности, обеспечив проверку всех графических узлов и механизма их генерации.

Такая архитектура позволила реализовать полноценное интеграционное тестирование. Поскольку на каждое действие пользователя в системе назначен свой обработчик событий~(callback), автотесты имитировали поведение человека, последовательно вызывая соответствующие функции-обработчики. В рамках этого этапа были реализованы сценарии сквозного тестирования: проверка компиляции базового графа, сборка модели простейшей линейной регрессии, а также конструирование сложной сверточной нейронной сети.

Проведенное тестирование подтвердило техническую надежность и стабильность разработанного MVP-конструктора. Автоматизированные и ручные проверки доказали, что спроектированная архитектура корректно справляется с построением и компиляцией графов любой сложности. Однако успешное прохождение программных тестов гарантирует лишь отсутствие багов в коде. Для оценки того, насколько созданный инструмент действительно снижает порог входа в разработку нейросетей и понятен неспециалистам, технической верификации недостаточно. Это обусловило необходимость следующего шага — проведения практической апробации продукта в реальных условиях с участием целевой аудитории.


\section{Практическая апробация конструктора в образовательной среде}

Для подтверждения практической значимости разработанного MVP-конструктора и оценки того, насколько успешно он решает проблему высокого порога входа в разработку искусственного интеллекта, была проведена масштабная апробация продукта в реальных условиях. Площадкой для внедрения стала Национальная технологическая олимпиада~(НТО) Джуниор - это одни из крупнейших в России командных инженерных соревнований, ориентированные на школьников 5–7 классов.

В рамках олимпиады существует несколько технологических направлений, одним из которых является сфера «Технологии и искусственный интеллект». В сезоне 2024/2025 года все соревнования в этой сфере были объединены сквозной научно-экологической легендой «Цифровой Байкал». Участники выступали в роли исследователей нового поколения, работающих на вымышленную компанию «Альтернативное будущее». Выбор данной площадки оказался идеальным для тестирования конструктора GraphNet: целевая аудитория~(учащиеся средней школы) не обладает профессиональными навыками программирования на Python и глубокими знаниями высшей математики, что полностью соответствует заявленной проблематике курсовой работы. Сами соревнования проходили в три последовательных этапа: отборочный этап, всероссийский финал и Слёт победителей.

На отборочном этапе участники работали индивидуально. Погружаясь в легенду, они анализировали экологическую обстановку Байкальского региона, а именно - изучали проблему природных пожаров. Перед школьниками стояли следующие задачи по работе с искусственным интеллектом.

\begin{enummarker}
\item Сбор и подготовка данных. парсинг информации о погоде и инфраструктуре поселений, формирование обучающей и тестовой выборок.
\item Проектирование нейронной сети. создание базовой архитектуры для задачи регрессии — предсказания количества пожаров. С помощью визуальных блоков участники собирали трехслойную полносвязную нейросеть.
\item Оценка качества модели. расчет метрики средней абсолютной ошибки~(MAE) с использованием специального вычислительного узла в конструкторе.
\end{enummarker}

Следует отметить, что довести это задание до готового решения смогло малое количество участников. Детальный анализ пользовательского опыта показал, что причиной стала не сложность самого графического инструмента. Задания по проектированию ИИ стояли в самом конце обширного списка олимпиадных задач~(после математики и фактчекинга), и у большинства участников  не хватило времени до них добраться в условиях жесткого ограничения времени отборочного тура.Участники, которые успели приступить к работе с конструктором, уверенно справлялись с логикой построения модели.

\myfigure{0.9}{image16.png}{Отборочный этап НТО Джуниор}{nto_stage1_img}

Финал олимпиады представлял собой интенсивное командное соревнование, посвященное тайне байкальской фауны и операции «Биоразнообразие». Работа строилась по принципу строгого распределения ролей: дата-журналист, дата-аналитик, ИИ-программист. Роль ИИ-программиста заключается в работе с GraphNet для создания системы компьютерного зрения «умного батискафа». На финале решались следующие задачи.

\begin{enummarker}
\item Создание архитектуры многоклассовой классификации. Основываясь на показателях датчиков алгоритм должен был классифицировать 5 видов организмов Байкала: голомянка, эпишура, байкальская губка и др. Для этого собиралась модель со специфической функцией активации на последнем слое: 35 selu - 15 selu - 5 softmax.
\item Компиляция и обучение: Настройка оптимизатора и функции потерь для правильной обработки нескольких классов.
\item Формирование предсказаний: прогон тестовых данных через обученную модель и сохранение результатов для загрузки на платформу проверки и интеграции в итоговое Вычислительное эссе команды.
\end{enummarker}

Как и на отборочном этапе, плотный график финала оставался серьезным барьером, но благодаря командному разделению труда большему числу участников удалось полноценно испытать конструктор на практике. Графовая структура позволила наглядно отследить, как данные от датчиков батискафа преобразуются в конкретные предсказания классов фауны.

\myfigure{0.9}{image17.png}{Финал НТО Джуниор}{nto_final_img}

Кульминацией апробации стал Слёт НТО Джуниор - очное образовательное мероприятие для победителей и призёров со всей страны, проходившее в формате свободного исследовательского хакатона. В отличие от предыдущих этапов с их стрессовыми временными рамками, здесь участники могли в спокойной обстановке погрузиться в настоящий Data Science.

\myfigure{0.9}{image18.png}{Слет НТО Джуниор}{nto_slet_img}

Им была предложена задача «на вырост», уровень которой выходит далеко за рамки школьной программы и соответствует реальным лабораторным работам студентов 2–3 курсов профильных IT-вузов. Школьникам предстояло работать со сложным набором данных CIFAR-10, состоящим из 60 000 цветных картинок~(размером 32x32 пикселя), разделенных на 10 классов~(самолеты, автомобили, птицы и т.д.). В рамках Слёта участники решали следующие продвинутые задачи.

\begin{enummarker}
\item Реализация классической архитектуры LeNet-5~\cite{bib_lenet_5} с нуля: самостоятельное проектирование цепочки обработки изображений. Школьники настраивали сверточные слои~(Conv2D) для извлечения линий и текстур из картинок, слои подвыборки~(Pooling) для уменьшения размера карт признаков, а также развертывание~(Flatten) для перевода найденных признаков в вектор и классификации через полносвязные слои~(Dense).
\item Адаптация модели под новые данные: перевод входного слоя с приема одного канала на три канала.
\item Модернизация и улучшение нейросети: проведение настоящих экспериментов по борьбе с переобучением. Участники увеличивали количество сверточных фильтров и заменяли устаревшие функции активации на современные ReLU и SeLU, чтобы сеть быстрее и эффективнее обучалась на цветном датасете.
\end{enummarker}

Показательно, что в условиях достаточного количества времени с реализацией и модернизацией сложной сверточной нейросети успешно справилась каждая команда.

Этот результат является исчерпывающим доказательством главной гипотезы исследования. Мощь и наглядность конструктора GraphNet позволили сделать сложную разработку нейросетей интуитивно понятной. Собирая архитектуру из визуальных блоков, пяти- и семиклассники сразу видели размерности тензоров на каждом слое и понимали, как данные проходят через нейросеть. Конструктор взял на себя всю рутинную нагрузку по написанию кода, позволив пользователям-новичкам полноценно применять системное мышление, фокусироваться на архитектуре искусственного интеллекта и решать реальные задачи машинного обучения.


\section*{Выводы по главе}
\addcontentsline{toc}{section}{Выводы по главе}
В ходе работы была спроектирована и реализована гибкая программная архитектура, основанная на иерархии типов аннотаций и механизме диспетчеризации, что позволило автоматизировать генерацию интерфейса узлов. Разработанный механизм событийного кэширования начальных узлов обеспечил высокую производительность компиляции графа вычислений.

Графическая оболочка приложения была реализована с учетом трех ключевых компонентов визуального языка программирования: синтаксиса~(цветовая индикация связей), семантики~(группировка узлов) и прагматики~(встроенная визуализация графиков обучения). Надежность системы подтверждена результатами многоуровневого тестирования, включая сквозные сценарии сборки сверточных нейронных сетей.

Практическая апробация продукта в рамках Национальной технологической олимпиады Джуниор продемонстрировала, что разработанный инструмент позволяет учащимся 5–7 классов успешно решать инженерные задачи университетского уровня, включая проектирование архитектуры LeNet-5. Результаты апробации полностью подтвердили гипотезу о том, что использование разработанного графического интерфейса значительно снижает порог входа в область искусственного интеллекта и избавляет пользователя от рутинного написания кода, сохраняя при этом широкие возможности для научных экспериментов.

\chapter*{ЗАКЛЮЧЕНИЕ}
\addcontentsline{toc}{chapter}{ЗАКЛЮЧЕНИЕ} 

В рамках данной курсовой работы была реализована поставленная цель — разработка MVP Desktop-приложения для графического конструирования нейронных сетей, позволяющего значительно снизить порог входа в область глубокого обучения.

В ходе выполнения работы были успешно решены все поставленные задачи.

\begin{enumarabic}
\item Проведен анализ существующих инструментов разработки нейросетей, который выявил разрыв между сложными профессиональными фреймворками и чрезмерно упрощенными сервисами типа «черный ящик». Это подтвердило необходимость создания гибкого, но интуитивно понятного визуального конструктора.
\item Разработана концепция графической среды, основанная на парадигме DataFlow. Были определены и реализованы три ключевых компонента визуального языка: синтаксис, семантика и прагматика.
\item Спроектирована архитектура программного комплекса, обеспечивающая высокую масштабируемость. Использование иерархии типов аннотаций и механизмов диспетчеризации позволило создать систему, в которую легко интегрировать новые слои и функции без изменения программного ядра.
\item Разработана программная основа проекта, включающая механизм компиляции графа на базе TensorFlow и графическую оболочку на DearPyGUI. Реализованный алгоритм событийного кэширования начальных узлов позволил оптимизировать процесс вычислений.
\item Проведена апробация продукта в реальных условиях на базе Национальной технологической олимпиады~(НТО) Джуниор.
\end{enumarabic}

Результаты апробации стали главным подтверждением практической значимости работы. Участники олимпиады — школьники 5–7 классов, не обладающие навыками программирования, — смогли в короткие сроки освоить интерфейс GraphNet и успешно решить задачи по созданию архитектур для многоклассовой классификации и регрессии. Особо значимым достижением стала работа участников на слёте победителей, где команды самостоятельно спроектировали и обучили сложную сверточную нейросеть LeNet-5 на датасете CIFAR-10~\cite{bib_cifar_10}.

Разработанный конструктор имеет потенциал для объединения школьного и вузовского образования. Он переводит фокус с рутинного написания кода на архитектурное мышление и понимание сути алгоритмов искусственного интеллекта.

Практическая значимость проекта заключается в создании готовой платформы, которая может быть использована в школьном и вузовском образовании, а также для быстрого прототипирования моделей специалистами нетехнических областей.

Научная и практическая новизна проекта заключается в создании гибкой архитектуры, где графическая оболочка динамически формируется на основе программных аннотаций, обеспечивая беспрецедентную расширяемость системы. В отличие от существующих образовательных аналогов, данный конструктор снимает «потолок сложности» и позволяет визуально проектировать нелинейные, ветвящиеся архитектуры любой глубины без написания кода. Успешная апробация продукта в рамках Национальной технологической олимпиады~\cite{bib_nto_junior}~(НТО) Джуниор подтвердила когнитивную эффективность выбранного подхода: учащиеся 5–7 классов смогли успешно решить инженерные задачи университетского уровня, включая проектирование сверточных сетей LeNet-5 для распознавания образов.

В дальнейшем планируется расширение функциональных возможностей системы, добавление генерации Python-кода из собранного графа и внедрение механизмов совместной работы над проектами в облачной среде.Оформи мою курсовую работу в этот LaTeX шаблон


\chapter*{СПИСОК ИСПОЛЬЗОВАННЫХ ИСТОЧНИКОВ}
\addcontentsline{toc}{chapter}{СПИСОК ИСПОЛЬЗОВАННЫХ ИСТОЧНИКОВ} 
\begin{thebibliography}{}

\bibitem{bib_goodfellow}
    Гудфеллоу, Я. Глубокое обучение / Я. Гудфеллоу, И. Бенджио, А. Курвилль ; перевод с английского А. А. Груздева. – 2-е изд., испр. – Москва : ДМК Пресс, 2018. – 652 с.

\bibitem{bib_vpl_theory}
    Shu, N. C. Visual Programming : [theory and concepts] / N. C. Shu. – New York : Van Nostrand Reinhold, 1988. – 315 p.

\bibitem{bib_orange}
    Demšar, J. Orange: Data Mining Toolbox in Python / J. Demšar, T. Curk [et al.] // Journal of Machine Learning Research. – 2013. – Vol. 14. – P. 2349–2353.

\bibitem{bib_tensorflow}
    Abadi, M. TensorFlow: A System for Large-Scale Machine Learning / M. Abadi, P. Barham [et al.] // 12th USENIX Symposium on Operating Systems Design and Implementation~(OSDI 16). – 2016. – P. 265–283.

\bibitem{bib_python_docs}
    Python 3.12 Documentation [Электронный ресурс]. – URL: https://docs.python.org/3/~(дата обращения: 24.03.2024).

\bibitem{bib_keras_api}
    Chollet, F. Deep Learning with Python / F. Chollet. – 2nd ed. – New York : Manning Publications, 2021. – 504 p.

\bibitem{bib_dearpygui}
    Dear PyGui Documentation: GPU Accelerated Python GUI Framework [Электронный ресурс]. – URL: https://dearpygui.readthedocs.io/~(дата обращения: 25.03.2024).

\bibitem{bib_gamma_patterns}
    Гамма, Э. Приемы объектно-ориентированного проектирования. Паттерны проектирования / Э. Гамма, Р. Хелм, Р. Джонсон, Дж. Влиссидес. – Санкт-Петербург : Питер, 2020. – 368 с.

\bibitem{bib_lenet_5}
    LeCun, Y. Gradient-based learning applied to document recognition / Y. LeCun, L. Bottou, Y. Bengio, P. Haffner // Proceedings of the IEEE. – 1998. – Vol. 86, no. 11. – P. 2278–324.

\bibitem{bib_cifar_10}
    Krizhevsky, A. Learning Multiple Layers of Features from Tiny Images / A. Krizhevsky. – Toronto : University of Toronto, 2009. – 7 p.

\bibitem{bib_nto_junior}
    Национальная технологическая олимпиада Junior: официальный сайт [Электронный ресурс]. – URL: https://ntojunior.ru/~(дата обращения: 26.03.2024).

\end{thebibliography}

\end{document}